% 第2章 2.3.1：保范数延拓的几何示意。
% (a) (R^2,‖·‖_1)：单位球为菱形，子空间 M 为横轴，ℓ(s,0)=s，‖ℓ‖=1。
%     延拓 L_t(x)=x_1+t x_2（|t|≤1）的对偶范数 max(1,|t|)=1，水平线 {L_t=1} 都过 (1,0) 且支撑菱形；
%     图中画 t=-1,-1/2,0,1/2,1。
% (b) (R^2,‖·‖_2)：单位圆在 (1,0) 只有一条支撑线 x_1=1，保范数延拓唯一。比例 1 单位=14mm。
\begin{tikzpicture}[x=14mm,y=14mm,font=\figurefont\footnotesize,text=FigureInk,
  note/.style={font=\FigureNoteFont\fontsize{7.5}{9}\selectfont,text=FigureMuted},
  ax/.style={draw=FigureStroke,line width=.6pt,-{Stealth[length=1.6mm]}}]
\begin{scope}
  \fill[FigureBlueFill] (1,0)--(0,1)--(-1,0)--(0,-1)--cycle;
  \draw[FigureBlue,line width=1.1pt] (1,0)--(0,1)--(-1,0)--(0,-1)--cycle;
  \draw[ax] (-1.4,0)--(1.75,0) node[below left=-.3mm] {$M$};
  \draw[ax] (0,-1.3)--(0,1.4);
  \foreach \t/\sty in {-1/FigureGreen,-0.5/FigureYellow,0/FigureInk,0.5/FigureYellow,1/FigureGreen}{
    \draw[\sty,line width=.8pt] ({1-\t*1.25},1.25)--({1+\t*1.25},-1.25);}
  \fill[FigureInk] (1,0) circle (1.3pt);
  \node[anchor=west] at (-1.5,1.65) {(a) $\ell^1$ 范数：延拓不唯一};
  \node[note,anchor=west] at (-1.5,1.18) {$x_1+tx_2=1$};
  \node[note,anchor=west] at (-1.5,0.95) {$|t|\leq1$};
\end{scope}
\begin{scope}[xshift=52mm]
  \fill[FigureBlueFill] (0,0) circle (1);
  \draw[FigureBlue,line width=1.1pt] (0,0) circle (1);
  \draw[ax] (-1.4,0)--(1.75,0) node[below left=-.3mm] {$M$};
  \draw[ax] (0,-1.3)--(0,1.4);
  \draw[FigureInk,line width=.8pt] (1,1.25)--(1,-1.25);
  \fill[FigureInk] (1,0) circle (1.3pt);
  \node[anchor=west] at (-1.5,1.65) {(b) $\ell^2$ 范数：延拓唯一};
  \node[note,anchor=west] at (1.08,0.9) {$x_1=1$};
\end{scope}
\end{tikzpicture}
